
% ----------------------------------------------------------------------
%                   LATEX TEMPLATE FOR PhD THESIS
% ----------------------------------------------------------------------

% based on Harish Bhanderi's PhD/MPhil template, then Uni Cambridge
% http://www-h.eng.cam.ac.uk/help/tpl/textprocessing/ThesisStyle/
% corrected and extended in 2007 by Jakob Suckale, then MPI-CBG PhD programme
% and made available through OpenWetWare.org - the free biology wiki


%: Style file for Latex
% Most style definitions are in the external file PhDthesisPSnPDF.
% In this template package, it can be found in ./Latex/Classes/
\documentclass[a4paper, openright,11pt]{Latex/Classes/PhDthesisPSnPDF}

%: Macro file for \LaTeX
% Macros help you summarise frequently repeated Latex commands.
% Here, they are placed in an external file /Latex/Macros/MacroFile1.tex
% An macro that you may use frequently is the figuremacro (see introduction.tex)
% This file contains macros that can be called up from connected TeX files
% It helps to summarise repeated code, e.g. figure insertion (see below).

% insert a centered figure with caption and description
% parameters 1:filename, 2:title, 3:description and label
\newcommand{\figuremacro}[3]{
	\begin{figure}[htbp]
		\centering
		\includegraphics[width=1\textwidth]{#1}
		\caption[#2]{\textbf{#2} - #3}
		\label{#1}
	\end{figure}
}

% insert a centered figure with caption and description AND WIDTH
% parameters 1:filename, 2:title, 3:description and label, 4: textwidth
% textwidth 1 means as text, 0.5 means half the width of the text
\newcommand{\figuremacroW}[4]{
	\begin{figure}[htbp]
		\centering
		\includegraphics[width=#4\textwidth]{#1}
		\caption[#2]{\textbf{#2} - #3}
		\label{#1}
	\end{figure}
}

% inserts a figure with wrapped around text; only suitable for NARROW figs
% o is for outside on a double paged document; others: l, r, i(inside)
% text and figure will each be half of the document width
% note: long captions often crash with adjacent content; take care
% in general: above 2 macro produce more reliable layout
\newcommand{\figuremacroN}[3]{
	\begin{wrapfigure}{o}{0.5\textwidth}
		\centering
		\includegraphics[width=0.48\textwidth]{#1}
		\caption[#2]{{\small\textbf{#2} - #3}}
		\label{#1}
	\end{wrapfigure}
}

% predefined commands by Harish
\newcommand{\PdfPsText}[2]{
  \ifpdf
     #1
  \else
     #2
  \fi
}

\newcommand{\IncludeGraphicsH}[3]{
  \PdfPsText{\includegraphics[height=#2]{#1}}{\includegraphics[bb = #3, height=#2]{#1}}
}

\newcommand{\IncludeGraphicsW}[3]{
  \PdfPsText{\includegraphics[width=#2]{#1}}{\includegraphics[bb = #3, width=#2]{#1}}
}

\newcommand{\InsertFig}[3]{
  \begin{figure}[!htbp]
    \begin{center}
      \leavevmode
      #1
      \caption{#2}
      \label{#3}
    \end{center}
  \end{figure}
}


%%% Local Variables: 
%%% mode: latex
%%% TeX-master: "~/Documents/LaTeX/CUEDThesisPSnPDF/thesis"
%%% End: 

\usepackage{url}	
\usepackage{amssymb,amsmath}
\usepackage{verbatim} 
\usepackage{amsmath}
\usepackage{array}
\usepackage{graphics}
\usepackage{subfigure}
\usepackage{polski}
\usepackage[utf8]{inputenc}

%: ----------------------------------------------------------------------
%:                  TITLE PAGE: name, degree,..
% ----------------------------------------------------------------------
% below is to generate the title page with crest and author name

%if output to PDF then put the following in PDF header
\ifpdf  
    \pdfinfo { /Title  (PhD and MPhil Thesis Classes)
               /Creator (TeX)
               /Producer (pdfTeX)
               /Author (YourName your@email.net)
               /CreationDate (D:YYYYMMDDhhmmss)  %format D:YYYYMMDDhhmmss
               /ModDate (D:YYYYMMDDhhmm)
               /Subject (xyz)
               /Keywords (add, your, keywords, here) }
    \pdfcatalog { /PageMode (/UseOutlines)
                  /OpenAction (fitbh)  }
\fi


%\title{Wykorzystanie środowiska PhysX do modelowania przepływów}
%
%
%
%% ----------------------------------------------------------------------
%% The section below defines www links/email for author and institutions
%% They will appear on the title page of the PDF and can be clicked
%\ifpdf
%  \author{\href{mailto:piotrjanisz@gmail.com}{Piotr Janisz}}
%%  \cityofbirth{born in XYZ} % uncomment this if your university requires this
%%  % If city of birth is required, also uncomment 2 sections in PhDthesisPSnPDF
%%  % Just search for the "city" and you'll find them.
%  \collegeordept{\href{http://www.something.net}{Faculty of Electrical Engineering, Automatics, Computer Science and Electronics}}
%  \university{\href{http://www.something.net}{AGH}}
%
%  % The crest is a graphics file of the logo of your research institution.
%  % Place it in ./0_frontmatter/figures and specify the width
%  \crest{\includegraphics[width=0.2\textwidth]{logo}}
%  
%% If you are not creating a PDF then use the following. The default is PDF.
%\else
%  \author{YourName}
%%  \cityofbirth{born in XYZ}
%  \collegeordept{CollegeOrDept}
%  \university{University}
%  \crest{\includegraphics[width=4cm]{logo}}
%\fi
%
%%\renewcommand{\submittedtext}{change the default text here if needed}
%\degree{M.Sc.}
%\degreedate{2012}


% ----------------------------------------------------------------------
       
% turn of those nasty overfull and underfull hboxes
\hbadness=10000
\hfuzz=50pt


%: --------------------------------------------------------------
%:                  FRONT MATTER: dedications, abstract,..
% --------------------------------------------------------------

\begin{document}

%\language{english}

% sets line spacing
\renewcommand\baselinestretch{1.2}
\baselineskip=18pt plus1pt


%: ----------------------- generate cover page ------------------------

%\maketitle  % command to print the title page with above variables
\thispagestyle{empty}

\begin{center}
Akademia Górniczo - Hutnicza im. Stanisława Staszica w Krakowie \\
Wydział Elektrotechniki, Automatyki, Informatyki i Elektroniki \\
Katedra Informatyki
% \\ Akademii Górniczo Hutniczej w Krakowie
\rule{\textwidth}{.1mm}
\end{center}

\begin{figure}[hkp!]
 \centering
% \includegraphics[width=50pt]{agh_znak}
 \includegraphics[bb=0 0 64 124]{logo}
 \label{fig:orzel_agh}
\end{figure}

% \vskip 80pt
\begin{center}
{\bf \large Tytuł pracy}
\vskip 10pt
{\bf \Large Wykorzystanie środowiska PhysX do modelowania przepływów}
\end{center}

\vskip 60pt
\begin{center}
{\Large \bf Piotr Janisz}\\
{piotrekjanisz@gmail.com}
\end{center}

\vskip 50pt
\begin{center}
{\bf
    Promotor pracy: dr inż. Witold Alda
}
\end{center}

\noindent
\rule{\textwidth}{.1mm}
\begin{center}
Kraków, 2012
\end{center}


%: ----------------------- abstract ------------------------

% Your institution may have specific regulations if you need an abstract and where it is to be placed in the document. The default here is just after title.


% Thesis Abstract -----------------------------------------------------


%\begin{abstractslong}    %uncommenting this line, gives a different abstract heading
\begin{abstracts}        %this creates the heading for the abstract page

Computing power growth allows more complex and accurate simulation techniques to be implemented in 
real-time applications. An example of those techniques are particle-based methods for simulation of fluids. 
They can provide a new level of realism into computer games. Nowadays most fluids in games are simulated using height field. 
Although using this method realistic waterbodies (like oceans or ponds) can be simulated, it's hard to achieve effects such as splashing or flooding.
Those can be easily simulated with particle-based methods such as Smoothed Particle Hydrodynamics (SPH). 

In this paper I would like to present realistic model of fluid that can be used in real time application, especially in computer games. Emphais will be placed on realistic rendering output from particle simulation.
For water simulation NVIDIA PhysX SDK will be used. For rendering particles I will use two different approaches: screen space rendering and isosurface extraction. Second algorithm is described in [reference] and is running on cpu. Authors presented performance of this algorithm running on one CPU core and I will present multithread implementation and it's performance analysis. At the end of this paper I will also present comparision of those two rendering techniques. 


\end{abstracts}
%\end{abstractlongs}


% ---------------------------------------------------------------------- 


% The original template provides and abstractseparate environment, if your institution requires them to be separate. I think it's easier to print the abstract from the complete thesis by restricting printing to the relevant page.
% \begin{abstractseparate}
%   
% Thesis Abstract -----------------------------------------------------


%\begin{abstractslong}    %uncommenting this line, gives a different abstract heading
\begin{abstracts}        %this creates the heading for the abstract page

Computing power growth allows more complex and accurate simulation techniques to be implemented in 
real-time applications. An example of those techniques are particle-based methods for simulation of fluids. 
They can provide a new level of realism into computer games. Nowadays most fluids in games are simulated using height field. 
Although using this method realistic waterbodies (like oceans or ponds) can be simulated, it's hard to achieve effects such as splashing or flooding.
Those can be easily simulated with particle-based methods such as Smoothed Particle Hydrodynamics (SPH). 

In this paper I would like to present realistic model of fluid that can be used in real time application, especially in computer games. Emphais will be placed on realistic rendering output from particle simulation.
For water simulation NVIDIA PhysX SDK will be used. For rendering particles I will use two different approaches: screen space rendering and isosurface extraction. Second algorithm is described in [reference] and is running on cpu. Authors presented performance of this algorithm running on one CPU core and I will present multithread implementation and it's performance analysis. At the end of this paper I will also present comparision of those two rendering techniques. 


\end{abstracts}
%\end{abstractlongs}


% ---------------------------------------------------------------------- 

% \end{abstractseparate}



%: ----------------------- contents ------------------------

\setcounter{secnumdepth}{3} % organisational level that receives a numbers
\setcounter{tocdepth}{3}    % print table of contents for level 3
\tableofcontents            % print the table of contents
% levels are: 0 - chapter, 1 - section, 2 - subsection, 3 - subsection



%: --------------------------------------------------------------
%:                  MAIN DOCUMENT SECTION
% --------------------------------------------------------------

% the main text starts here with the introduction, 1st chapter,...
\mainmatter

\renewcommand{\chaptername}{} % uncomment to print only "1" not "Chapter 1"


%: ----------------------- subdocuments ------------------------

% Parts of the thesis are included below. Rename the files as required.
% But take care that the paths match. You can also change the order of appearance by moving the include commands.


% this file is called up by thesis.tex
% content in this file will be fed into the main document

%: ----------------------- introduction file header -----------------------
\chapter{Introduction}

\graphicspath{{introduction/figures/}}

% ----------------------------------------------------------------------
%: ----------------------- introduction content ----------------------- 
% ----------------------------------------------------------------------

\section{Motivation}

\section{State of the art}

% ----------------------------------------------------------------------



	% background information


% this file is called up by thesis.tex
% content in this file will be fed into the main document

%: ----------------------- introduction file header -----------------------
\chapter{SPH}


\graphicspath{{sph/figures/}}

This chapter will provide brief introduction to modeling method used in my project - Smoothed Particle Hydrodynamics (SPH).

\section{Introduction}
To perform numerical simulation of natural phenomena several steps have to be performed (see \cite[section 1.1.2]{Liu}). First the mathematical model has to be obtained from observation of physical phenomena. The mathematical model consists of a set of governing equations, boundary conditions and initial conditions. Next governing equations have to be discretized in order to solve them with a numerical algorithm.
For describing physical governing equations two basic techniques exists: Eulerian description and Lagrarian description.

For describing governing equations of fluid flows two basic techniques exists: Lagrarian and Eulerian. In the first one we look at the fluid as a set of parcels which moves though space and time. Each fluid parcel has it's own mass which is fixed, and is identified by it's position $x_0$ at time $t = 0$. We also assume that we can obtain a function of fluid parcels position $x = x(t, x_0)$. This is demonstrated on figure~\ref{eulerian_lagrarian}.
Eulerian description assumes that we observe fluid that moves through fixed spatial points. Thus we are given velocity field $v$ at every spatial point $x$ and time $t$ - $v = v(x, t)$ (figure~\cite{eulerian_lagrarian}). 
More on fluid field description can be found in \cite[section 2.1]{Hauke2008} and \cite{Price06}

\figuremacro{eulerian_lagrarian}{Different fluid observation techniques}{on the left we can see Eulerian velocity field, on the right Lagrarian pathlines.}

For discretization there are several approaches: grid based (lagrarian grid and eulerian grid) and meshfree methods (SPH, particle in cell).
Lagrarian grid is attached to material and moves with it. Each mesh cell represents one fluid parcel from Lagrarian description. The advantages of this method are simpler equations given by lagrarian description, less memory usage (as the mesh is specified only where fluid or other material is present), easiness of tracking movement of fluid parcels. However this method does not handle large material deformations over time (like for instance water splashes) well. In such cases mesh must be recomputed and it an expensive procedure.
Eulerian grid, on the other hand, is fixed in space, and covers entire domain in which fluid can move. In each grid node we calculate mass momentum and energy flux. This method handles large deformations of material easily but its computional complexity is large, especially in three dimensions. The amount of memory require is also large as whole domain must be covered with grid. 
Third way of discretization are mesh free method. The idea of them is to combine strong points of both eulerian and lagrarian grids, but without using a grid or a mesh. This is done usually by discretizing material as a set of particles. At each particle values of functions, derivatives and integrals are approximated with particle approximations. An example of mesh free method is SPH which will be further described in this chapter. 
Lagrarian and Eulerian grids as well as meshfree methods are described in more details in \cite[chapter 1]{Liu}

\section{Fundamentals of SPH}

The main idea of SPH is to replace given field A with it's integral interpolant (or kernel interpolant):
\begin{equation}
\label{eq:sph_kernel_interpolant}
A_I(r) = \int A(r')W(r - r', h)dr'
\end{equation}

Where W is a smoothing kernel and has following properties:

\begin{equation}
\label{eq:sph_kernel_property_1}
\int W(r - r', h)dr' = 1
\end{equation}

\begin{equation}
\label{eq:sph_kernel_property_2}
\lim_{h \to 0}W(r - r', h) = \delta(r-r')
\end{equation}

\begin{equation}
\label{eq:compact_condition}
(\exists \kappa \in \mathbb{R}_+ )(W(r-r', h) = 0~when~|r-r'| > \kappa h)
\end{equation}

Integral interpolant \ref{eq:sph_kernel_interpolant} can be further approximated by sumation interpolant:

\begin{equation}
\label{eq:sumation_interpolant}
A_s(r) = \sum_{b}m_b\frac{A_b}{\rho_b}W(r-r_b,h)
\end{equation}

where $\frac{m_b}{\rho_b}$ is a volume of particle $b$. What SPH does is it divides space into set of particles, each particle has it's own mass and density. Mass is fixed and density can change thus changing particle's volume. When we want to compute value of A in some point in space $r$ we compute weighted average of A at neighboring particles. Weight depends on distance of particle's center form $r$ and on particles volume. 
It's also worth noting that \ref{eq:compact_condition} defines finite domain in which smoothing kernel is non-zero. This allows to speed up computations by summing over particles in nearest neighborhood.

Replacing field $A$ with it's kernel interpolation $A_I$ yields error of $h^2$ (see~\cite[section 2.2.1]{Liu}). On the other hand accuracy of sumation interpolant \ref{eq:sumation_interpolant} is hard to predict (see~\cite[section 12.1]{Monaghan1992}).

In many physical simulations we need to compute gradient of some field A. It turns out that using integral interpolant \ref{eq:sph_kernel_interpolant} we only need to compute gradient of smoothing kernel $W$ (\cite[section 2.2.2]{Liu}):

\begin{equation}
\label{eq:integral_interpolant_gradient}
\triangledown A_I(r) = \int A(r') \triangledown W(r - r', h)dr'
\end{equation}

As earlier we can aproximate integral with sumation:

\begin{equation}
\label{eq:sumation_interpolant_gradient}
\triangledown A_s(r) = \sum_{b}m_b\frac{A_b}{\rho_b}\triangledown W(r-r_b,h) 
\end{equation}

SPH is a interpolation method - does not uses a grid (MESHFREE), approximates fluid or other materials as a set of particles. Requires less resources comparing to grid methods when it comes to such phenomenas as splasches or breaking waves. 

SPH first used by Leon Lucy and R. A. Gingold and J. J. Monaghan

Different kernels can be used for different computations.


paper on astrophysics from 1992 (Monaghan1992) - \cite{Monaghan1992}

book about sph (Liu) - \cite{Liu}

Particle based simulations for interactive applications (Muller2003) - \cite{Muller2003}

\section{SPH for fluids}
Fluid motion can be described by two equations. First one is conservation of mass:
\begin{equation}
\label{eq:mass_conservation}
\frac{DM}{Dt} = 0
\end{equation}
second one is Navier-Stokes conservation of momentum equation:
\begin{equation}
\label{eq:momentum_conservation}
\rho \frac{Dv}{Dt} = -\triangledown p + \mu \triangledown^2v + f
\end{equation}

Both equations are presented in Lagrarian form. $\frac{D}{Dt}$ is a material derivative (TODO opisać co to jest, referencja).  $p$ is the pressure, $\mu$ is fluid viscosity and f represents external forces acting on the fluid. 

Navier-Stokes equations: 
 - continuity, 
 - momentum conservation, 
 - energy conservation
For SPH first one is not necessary and for fluid simulation in games energy conservation equation can be discarded as we only interested in fluid movement, not temperature (we can assume the temperature is constant). 

TODO SPH equations needs to be symmetrized.




% this file is called up by thesis.tex
% content in this file will be fed into the main document

%: ----------------------- introduction file header -----------------------
\chapter{Tools}


\graphicspath{{tools/figures/}}

% ----------------------------------------------------------------------
%: ----------------------- introduction content ----------------------- 
% ----------------------------------------------------------------------

This chapter will present brief introduction of tools used in building visualization system. This is no intended to be a tutorial as those information can be found on the Internet. 

\section{OpenGL}
OpenGL is an API specification for interacting with graphic hardware in order to produce 2D and 3D graphic. 
As a specification it is not bound to any programming language, operating system or hardware. 
OpenGL bindings exists for many languages including C/C++, Java, Python, Ruby and C$\sharp$ (see~\cite{OglBindings}). 

Since introduction of version 3.2, specification has been divided into core profile and compatibility profile. Core profile is smaller and contains only modern style API introduced with version 3.0. Compatibility profile implements all legacy functionality prior to OpenGL 3.0, including fixed pipeline. 

Fixed function pipeline is built-in to hardware and there is no control over how operations are performed on GPU. That means operations like vertex transformation, algorithms for shading surfaces are already implemented and can't be changed. The behavior may only be changed by choosing one of predefined methods or adjusting available parameters. Pros are that it's faster and less error prone to build an application. Programmers only have to choose existing technique and specify geometry primitives. On the other hand when new algorithm is invented there is no possibility to use it. There are also no possibilities to optimize existing techniques for the particular application. 
That is why programmable pipeline was introduced in modern GPUs. Rendering stages can be customized by writing programs called shaders. At first those programs were very specific and allowed little customization. Over time more and more capabilities were added to shaders and now their usage goes beyond graphic rendering \cite[Part~VI]{GPUGems3}. Currently programmers have full control over how primitives are rendered on the screen.

Apart from core and legacy profiles there are extensions. This is a mechanism of including vendor specific functions, implemented only on some GPUs, that can make use of some new hardware features. 

OpenGL is widely available on different hardware from graphic cluster, through desktop computers, game consoles and portable devices. For the last group dedicated standard has been created - OpenGL ES. It a subset of regular specification designed to run efficiently on devices with limited hardware resources. The most frequently used and the most useful functionality of OpenGL was included in OpenGL ES. That was desired as larger API requires more complex and larger driver to support it. It's worth mentioning that OpenGL ES is the only officially supported 3D API for portable devices running under Android and iOS operating systems. 

\figuremacro{opengl_pipeline}{OpenGL 3.3 pipeline}{taken from \cite[chapter~12]{OpenGLSuperbible}}

Figure~\ref{opengl_pipeline} shows diagram of OpenGL 3.3 pipeline. Three stages are fully programmable - vertex shader, geometry shader and fragment shader and those will be described in more details in following subsections. Also some of more important features of OpenGL used in implementing visualization algorithms will be described.

\subsection{Vertex Shader}
Input to the vertex shader program is a single vertex (vertices are processed in parallel) with optional attributes (which can vary between vertices - i.e.  normal vector, color, texture coordinate) and uniforms (which have the same value for all vertices - i.e. projection matrix, model matrix). This stage typically transforms vertices from model space to view space and applies perspective projection. Other typical applications are calculating perspective lighting and texturing calculations. More advanced operations includes procedural animation by modifying the position of the vertex. This can be used to animate water surfaces (like pond or oceans), clothes or skin. On modern GPUs the vertex shader has also access to texture data - especially useful when using textures as data-structures in physical simulations (see \cite[pages~412-419]{OpenGLSuperbible}).

\subsection{Geometry Shader}
This stage takes entire primitives (triangles, lines, points) as an input. The geometry shader can produce new primitives, discard existing or change their type. Thus it is capable of changing amount of data in pipeline, which is unique in contrast to vertex and fragment shaders (fragment shaders can only discard fragments). There are many examples of using geometry shader - from simple as rendering six faces of a cube map \cite{Gregory2009}, drawing vertices normals \cite[pages~434-437]{OpenGLSuperbible} to more advanced as shadow volume extrusion \cite[section~10.3.3.1]{Gregory2009}, isosurface extraction \cite{3DCourseSIGGraph2007} and dynamic tessellation. 

\subsection{Fragment Shader}
This stage operates on fragments. One fragment is usually one pixel on the screen, except cases when multisampling or supersamplings are used - then several fragments may contribute to one pixel. The fragment shader's job is to process pixels by computing per pixel lighting or applying textures. Input to this stage are per-fragment attributes passed from vertex or geometry shaders, uniform attributes which are passed before rendering and which have the same value for all fragments and textures. As vertex and geometry shaders operates on vertices, per-fragment attributes passed to fragment shader have to be interpolated. There are two types of interpolation flat and smooth. Flat means that all fragments in the primitive has the same value of attribute. Smooth means that values are interpolated between edges of primitive. 

As it operates at pixel level it is perfect for post processing effects (the one added on rendered scene) like blurring, bloom, different sort of filtering. In my project I make heavy use of fragment shaders. 

\subsection{Perspective projection and homogeneous space}
Perspective projection mimics the kind of image produced by a typical camera by providing effect known as perspective foreshortening. With this effect the farther object lies from camera the smaller it appears on the screen. 
Perspective projection transforms points form view space into homogeneous clip space. In OpenGL view space uses right handed coordinate system (figure~\ref{perspective_projection}) and homogeneous space uses left handed coordinate system (figure~\ref{homogenous_space}). 
The region that is visible from camera is called the view volume or view frustum (figure~\ref{perspective_projection}). It is defined by six  planes - near plane, far plane and four side planes. Only objects inside view volume appears on the screen after rendering. Objects are projected on near plane, part of which lies inside four side planes is a virtual screen.
Virtual screen left, right, bottom and top edges lies at $x=l, x=r, y=b, and y=t$. It's z coordinate is equal to -n. Those parameters are used to computed projection matrix $M_{V \to H }$ (equation~\ref{eq:projection_matrix}) ($V \to H$ means transformation from view space into homogeneous space):
\begin{equation}
\label{eq:projection_matrix}
M_{V \to H } = 
\begin{bmatrix}
\frac{2n}{w} & 0            & 0                & 0                \\ 
0            & \frac{2h}{n} &                  & 0                 \\ 
0            & 0            & -\frac{f+n}{f-n} & -\frac{2fn}{f-n} \\ 
0            & 0            & -1               & 0
\end{bmatrix}
\end{equation}
where $w=r-l$ and $h=t-b$.

\myFigureMacroW{perspective_projection}{A perspective view volume (frustum)}{1.0}

\myFigureMacroW{homogenous_space}{The canonical view volume in homogeneous clip space}{0.7}



\subsection{Depth buffer} \label{sec:depth_buffer}
Dept buffer (also called z-buffer) is used to determine which fragments occlude other fragments - thus which should be drawn on the screen. This is done by comparing values of z component of fragments from homogeneous space. 
After multiplying point from view space - $p_V = [x, y, z, 1]$,  by $M_{V \to H}$ we get point $p_H$:
\begin{equation}
\label{eg:projection_transformation}
p_H =
\begin{bmatrix}
\frac{2n}{w} & 0            & 0                & 0                \\ 
0            & \frac{2h}{n} &                  & 0                 \\ 
0            & 0            & -\frac{f+n}{f-n} & -\frac{2fn}{f-n} \\ 
0            & 0            & -1               & 0
\end{bmatrix}
*[x, y, z, 1]^T
\end{equation}
thus we get:
\begin{equation}
\label{eg:after_perspective_projection}
p_H = [\frac{2n}{w}x, \frac{2n}{h}y, -\frac{f+n}{f-n}z - \frac{2fn}{f-n}, -z]
\end{equation}
To transform $p_H$ into 3D space we need to divide x, y and z components by $\omega$ component, which is equal to $-z$. That gives us:
\begin{equation}
\label{eg:after_perspective_projection}
p_H = [-\frac{2nx}{wz},-\frac{2ny}{hz}, \frac{f+n}{f-n} + \frac{2fn}{z(f-n)}]
\end{equation}
After transformation, z component of points inside view frustum varies from [-1, 1]. This has to be transform to [0, 1], thus value of depth buffer d(z) as follows:
\begin{equation}
\label{eq:depth_value}
d(z) = \frac{\frac{f+n}{f-n} + \frac{2fn}{z(f-n)} + 1}{2}
\end{equation}
Equation~\ref{eq:depth_value} shows that value of depth buffer is a nonlinear function of z component of the point. Figure~\ref{z_buffer_dist} shows how this value changes with changing near plane distance from the viewer. It can also be seen that precision is not evenly distributed and it decreases with the distance from near plane. This can lead to so called z-fighting which occurs when two objects have close z-buffer values. When the objects are close to the viewer z-values are distinguishable, however when objects get further, values can become equal due to lack of precision (see~\ref{z_fighting}).

Z-fighting can be avoided when reasonable values are chosen for far and near planes. However there are cases when evenly distributed values are required (see section~\ref{sec:gaussiansmoothing}). Alternative technique to Z-buffering is W-buffering which offers better distribution of depth values \cite{Gregory2009}. Some hardware supports W-buffering but most of it doesn't and it doesn't seem to be supported in future. A workaround to this problem was presented in \cite{Dunlop2006}. It customizes projection transformation in vertex shader in a way that depth buffer values have linear distribution at the end. 

Linear depth value can also be obtained from equation~\ref{eg:after_perspective_projection} by dividing $\omega$ component of vector $p_H$ by $f$:
\begin{equation}
\label{eq:linear_depth}
d_{lin}(z) = \frac{-z}{f}
\end{equation}
since $-z$ is in range [n, f] value of $d_{lin}(z)$ will be in range $[\frac{n}{f}, 1]$. 

\figuremacroW{z_fighting}{Z-fighting}{two polygons rendered - on the left when z-fighting occurs, on the right the same polygons rendered correctly.}{1.0}
\myFigureMacroW{z_buffer_dist}{Value of z buffer as a function of distance from near plane for different values of n. $f-n$ is kept constant at 100}{1.0}

\subsection{Framebuffer Object}
Normally after performing drawing operations, output is directed to framebuffers provided by windowing system. From those buffers data can be displayed on the computer screen. However, sometimes it's not desired for data to be displayed on the screen but to be used as an input to the rendering pipeline again. This technique is called off screen rendering. It can be achieved by copying data from default framebuffers to textures and than used as an input to the pipeline again but this requires additional copy operation that should be avoided. Framebuffer objects are mechanism for dealing with situation when data shouldn't be send directly to the display. Frame buffer objects are containers that can hold other object that have memory storage and can be rendered to \cite[chapter~8]{OpenGLSuperbible}. Two kinds of objects that can be rendered to are renderbuffer objects and textures. Rendering to textures is particularly useful as data doesn't have to be copied to be used as an input to the pipeline. This allows to efficiently implement rendering techniques that requires multiple phases.
Unlike default framebuffer object, the one created by user is not limited in size to the size of rendering window. This can be used to speedup some computationally expensive processing by rendering with lower resolution into texture and than applying this texture to the original window with linear filtering. 

Details about framebuffer objects, their creation and attaching storage to them can be found in \cite[chapter~8]{OpenGLSuperbible} and \cite[chapter~10]{RedBook}.

\section{Boost library}
Boost is a set of libraries extending capabilities of C++ language. Libraries provides classes and functions for most common task and algorithms, like regular expressions, hash maps, threading. They are also cross platform. 

The advantage of boost libraries is that they are designed for maximum flexibility and speed. They make heavy use of C++ template programming to be as general as possible. The other advantage of boost libraries is that they often become included in C++ standard library - like regular expressions (see~\cite{CppRef}) . 

On the other hand relying on templates makes compiling times longer. Template classes and functions must be defined entirely in header files what makes it impossible to compile them into object files. Thus whenever file has to be recompiled template classes have to be generated again.

Boost libraries used in my project are \textit{boost regex} and \textit{boost threads}. The first one is used to process configuration files. \textit{Boost threads} library is used to parallelize isosurface extraction algorithm presented in section~\ref{sec:isosurfaceextraction}.

\section{PhysX}

PhysX is a framework for performing physical simulations in real time. It's designed to be used in computer games and is optimized for performance. It provides most necessary models required in games: rigid bodies, soft bodies, cloths, fluids and joints. What is more it allows simulations to be performed on GPU which is much faster especially when large amounts of object acts in a simulation.

The main drawback is that hardware acceleration can only be performed on Nvidia's GPUs. 

PhysX works in asynchronous way - that means simulation step is computed either on GPU or CPU thread and other CPU cores can perform some other work in the meantime - for example render scene using data from previous step. The asynchronous workflow of PhysX is presented on figure~\ref{physx_workflow}.

\myFigureMacroW{physx_workflow}{Diagram of PhysX SDK workflow}{0.7}

My project uses two PhysX features - rigid body and water simulations and thus those will be described in more details. 

% --------------------------------- Rigid bodies 
\subsection{Rigid bodies}

Basic entity of simulation is an actor. Actors can be either static objects, fixed in space (walls, ground, etc.) or dynamic rigid bodies.
Figure~\ref{physx_actors} shows different ways the same actors can be represented. Actors have shapes assigned to them which are used by collision detection system. This shape can be the same mesh as used by rendering system, however it is much more efficient to approximate it by set of bounding volumes (figure~\ref{physx_actors}). Bounding volumes are simple shapes for which collision detection is easy and fast - usually rectangles and spheres. 

For physics computation, actor is represented as a tensor located at the center of mass (figure~\ref{physx_actors}). In case of collision with other object contact points are provided as points in space with normal vector of collision planes (red arrows on figure~\ref{physx_actors}). Center of mass is automatically computed from actor's shape, however there is a possibility to set it manually.
\figuremacroW{physx_actors}{Different ways of representing actors}{From left: graphical representation for rendering engine, bounding volumes used by collision detection and finally representation used for dynamics.}{1.0}

Important property of actor is its frame. It defines actor local coordinates system. Actor's frame is given as a matrix which transforms points from actor's local coordinates into world coordinates. All actor's shapes positions and orientations are given relatively to actor's frame.

Body movement can be a combination of rectilinear (along path) and angular (around center of mass) motions. Dynamic actors have set of properties used for rigid body simulation. All properties have a linear and angular version. They are gathered in table~\ref{tab:rigid_body_properties}

\begin{table}[htdp]
\caption[Rigid body parameters]{\textbf{Comparision of sphere rendering methods}}
\centering
    \begin{tabular}{|p{6.5cm}|p{7cm}|}
        \hline
        {\bf Linear property } & {\bf Angular property } \\ \hline
       {\bf  mass }          & {\bf  interia } - mass distribution of the body along its three dimensions         \\ 
        {\bf  position } - center of mass position of rigid body       & {\bf  orientation } - 3x3 matrix representing  principal moments relative to the actor's frame     \\ 
        {\bf  velocity } - linear velocity vector      & {\bf  angular velocity } - vector representing axis of rotation which length is equal to magnitude of angular velocity \\ 
        {\bf  force } - vector representing force acting on body     & {\bf  torque } - angular force acting on a body represented as axis of rotation with length equal to magnitude of angular velocity          \\
        \hline
    \end{tabular}
\label{tab:rigid_body_properties}
\end{table}

Physical properties presented in table~\ref{tab:rigid_body_properties} controls how body reacts when force or torque is applied. Actors behavior after collision is determined by theirs material properties. Those parameters are: restitution, static friction and dynamic friction. Normally material is applied to entire actor's shape, however there is a possibility to apply materials per triangle. Additionally friction parameters can be anisotropic. 
% --------------------------------- Fluids
\subsection{Fluids}
PhysX provides three methods for simulating fluids: SPH, simple method and mixed method. First one was described in chapter~\ref{chap:sph}. Second one does not take interparticle forces into consideration and thus simplifies computations. This simplified method is good for simulating phenomenas such as smoke or fog, but for realistic water SPH must be used. Third one - mixed method, according to documentation \cite{PhysXDoc}, alternates between SPH and simple mode, providing more performance than SPH but still maintaining some dense characteristics.

Implementation details of PhysX's SPH method are covered in \cite{Harris08}. It solves conservation of momentum equation~\ref{eq:momentum_conservation} to move particles. External force includes only gravity. Symmetrization of pressure and viscosity forces are done as in equations \ref{eq:sph_pressure_force_muller} and \ref{eq:sph_viscosity_force_sym}. Surface tension (see~\ref{eq:surface_tension}) is also modeled, however there are no implementation details given on that.

There are two ways of creating fluids in PhysX - either directly specifying particles initial positions, velocities, etc. or by using emitters (figure~\ref{emitters_drains}). Emitters are objects that generates particles. They can be static objects or they may be attached to a rigid body and interact with other objects in the scene. Emitters can generate particles in two modes:
\begin{itemize}\itemsep2pt
\item \textit{Constant pressure} - in which state of surrounding fluid is taken into account. Emitter tries to match rest spacing of particles (which is one of fluid's parameters). 
\item \textit{Constant flow} - emitter keeps generating the same number of particles each frame. 
\end{itemize}

\figuremacro{emitters_drains}{PhysX emitters and drains}{Simple 2D scene with one emitter and two drains.}

Removing particles can be done in three ways:
\begin{itemize}\itemsep2pt
\item By removing particles manually.
\item By setting particles lifetime parameter which is time in seconds an emitted particle lives.
\item By creating drain objects (figure~\ref{emitters_drains}). When particle hits drain it is removed from simulation. Those object are useful to prevent fluid from spreading too far and negatively affecting performance.
\end{itemize}

Fluids can interact with other objects in the scene. The interaction is handled bypassing conservation of momentum equation~\ref{eq:momentum_conservation} by applying velocity changes directly to particles. Interaction with rigid bodies can be either one-way (rigid body acts on fluid) or two-way (when fluid also acts on rigid body).

Most important fluid parameters that can be set are:
\begin{description}\itemsep1pt
\item[Kernel Radius Multiplier] Corresponds to $k$ parameter of SPH smoothing kernel. Together with \textit{Rest Particle Per Meter} defines particle's radius of influence (area within smoothing kernel is non-zero): $radius = \frac{kernelRadiusMultiplier}{restParticlesPerMeter}$
\item[Rest Particles Per Meter] Defines number of particles per volume unit when fluid is in rest state.
\item[Rest Density] Corresponds to $\rho_0$ from equation~\ref{eq:pressure}. Defines density of fluid in it's rest state.
\item[Viscosity] Corresponds to $\mu$ constant in equation~\ref{eq:momentum_conservation}. Lower viscosity will produce more runny fluids like water, higher viscosity will produce more sticky fluids like honey.
\item[Stiffness] Corresponds to $k$ constant in equation~\ref{eq:pressure}. Defines how compressible the fluid is. Higher values means less compressible fluids (for instance water). 
\item[Surface Tension] Equivalent of $\sigma$ from equation~\ref{eq:surface_tension}. The higher this parameter is the smoother fluid surface will be. 
\end{description}

PhysX allows creating different fluids with different parameters and simulating them in the same time. However those fluids can't interact. That means it is not possible to simulate mixing two fluids together. Another limitation of PhysX's fluids is limit of 64000 particles per fluid.
% ----------------------------------------------------------------------





% this file is called up by thesis.tex
% content in this file will be fed into the main document

\chapter{Rendering techniques} % top level followed by section, subsection

\graphicspath{{2/figures/}}

% ----------------------- contents from here ------------------------

Output from particle base simulations is list of particles containing positions. It can also contain additional parameters - for example PhysX fluid simulation returns velocity, lifetime and density in addition to position for each particle. 
Traditionally triangle meshes are used for rendering objects. Using this approach requires constructing fluid triangle mesh for given particle positions. This method is called isosurface extraction and will be described in second section of this chapter. 
Another possibility is to render each particle as a point (or a billboard in general) with opacity so that when large amount of particles is rendered the result would give illusion of a smooth surface. Such technique has one major drawback - it does not produces surface prevents us from adding effects of reflection and refraction.
Similar technique is to extract visible surface by rendering each particle as a sphere into depth buffer. This will be described in first section.

\section{Screen space}
\figuremacroW{screen_space}{Screen Space Fluid Rendering}{The figure shows high level overview of the method, taken from \cite{Green2010}}{0.7}
As mentioned before this approach extracts visible surface by rendering each particle into depth buffer. High level overview ot this method is presented on figure~\ref{screen_space} and consists of following steps \cite{laanSainz2009}: rendering depth texture (section~\ref{sec:surfacedepth}) and thickness textures (section~\ref{sec:thickness}), depth smoothing (section~\ref{sec:smoothing}) and assembling fluid surface with rest of the scene (section~\ref{sec:rendering}).

\subsection{Surface depth} \label{sec:surfacedepth}
To obtain fluid surface visible from the viewpoint of camera each particle is rendered as a sphere into depth texture [rysunek]. At each pixel only closest value is kept using hardware depth test. 
To avoid rendering large amount of geometry each particle is rendered as a point sprite and it's depth is generated in fragment shader. This common technique speeds up rendering process significantly as well as improves quality of rendered spheres (as can be seen on figure~\ref{spheres}) because depth values are generated for each pixel. Rendering spheres as point sprites is 10 to 100 times faster comparing to rendering them as triangle meshes (see table~\ref{sphere_speed_comp}).
\figuremacroW{spheres}{Spheres rendered with different methods}{1 - mesh with 128 triangles, 2 - mesh with 512 triangles, 3 - mesh with 2048 triangles, 4 - point sprite with normals generated in fragment shader}{0.5}
\begin{table}[htdp]
\caption[Comparision of sphere rendering methods]{\textbf{Comparision of sphere rendering methods}}
\centering
\begin{tabular}{cr} 
{\bf Method} & {\bf Frames per second} \\ 
\hline 
Mesh, 128 triangles & 20 \\
Mesh, 512 triangles & 6 \\
Mesh, 2048 triangles   & 1 \\
Point Sprites & 200 \\
\end{tabular}
\label{sphere_speed_comp}
\end{table}

\subsection{Smoothing} \label{sec:smoothing}
Although previous step produces surface it's quality is not sufficient. As can be seen on figure~\ref{spheres_no_smoothing} individual spheres can be seen giving fluid surface gelly-like appearance. To remove this artifact some kind of smoothing has to be applied. Section~\ref{sec:gaussiansmoothing} will describe gaussian smoothing and section~\ref{sec:curvatureflowsmoothing} curvature flow smoothing.
\figuremacroW{spheres_no_smoothing}{Fluid surface rendered without smoothing phase}{Individual spheres can be seen, giving surface unnatural appearance}{0.7}

\subsubsection{Gaussian smoothing} \label{sec:gaussiansmoothing}
Most obvious way to smooth values in depth texture is to apply gaussian filter. It's easy to implement and can be computed fast due to it's linear separability. However this filter produces undesired effect of blending drops of fluid with background surfaces (see figure~\ref{gaussian_filter_img}). 
\figuremacroW{gaussian_filter_img}{Fluid surface rendered with Gaussian smoothing}{From left: depth image, depth image after applying Gaussian filter, diffuse shaded surface}{0.7}
Thus edge-preserving filters needs to be used which are also called bilateral filters. Bilateral Gaussian filter is a modification that changes wages of pixels depending on difference between their tonal value ($I(s)$) and tonal value of central pixel ($I(s_0)$). This can be described by following formula (from~\cite{PhamVliet2005}): 
\begin{equation}
\label{bilateral_equation}
O(s_0) = \frac{\sum_{s \in S}f(s, s_0)I(s)}{\sum_{s \in S}f(s, s_0)}
\end{equation}
where 
\begin{equation}
\label{bilateral_weights_equation}
f(s, s_0) = g_s(s-s_0)g_t(I(s)-I(s_0))
\end{equation}
is the bilateral filter for the neighborhood around $s_0$. $g_s$ is a spatial weight, $g_t$ is a tonal weight and they both are Gaussian functions:
\begin{equation}
\label{bilateral_dist_equation}
g_s(s) = g(x, \sigma_S)g(y, \sigma_S)  \quad   g_t(I) = g(I, \sigma_t)
\end{equation}
As can be seen on equation~\ref{bilateral_weights_equation} only change in bilateral filter (in comparison to regular bilateral filter) is introduction of tonal weight $g_t$. This change makes filter space-variant - that means it can't be computed as a product of two one dimensional filters ($g_s$ in equation~\ref{bilateral_dist_equation}). 
Computional complexity of space-variant filters is $O(Nm^d)$ comparing to only $O(Nmd)$ for space-invariant (d is image dimensionality, m is the size of filtering kernel and N is the number of pixels in the image). For smoothing fluid surface kernel with $m = 20$ must be used, which gives for 2 dimensional image about 400 operations for each pixel when using bilateral filtering. In comparison separable implementation requires 40 operations. It turns out however that computing bilateral filter as it was a space-invariant gives good approximation for real time applications (see figure~\ref{bilateral_comp}). Approximation gives some artifacts but they are not visible when fluid is moving and other effects (reflection and refraction) are applied.
\figuremacroW{bilateral_comp}{Comparision of normal bilateral filter with it's separable approximation}{Upper row presents images for bilateral filter and bottom row presents separable approximation.}{0.7}

Bilateral filter has some undesired effect when applied on z-buffer output. The problem is that depth values are not evenly distributed. This mean that difference between two fragments depth values is dependent on their distance to the viewer. As can be seen on equation~\ref{bilateral_weights_equation} bilateral weights takes into account difference between depth values. In result the further surface lies from camera the less edges are preserved (see figure~\ref{edge_preservation}). Alternative technique to Z-buffering is W-buffering which offers better distribution of depth values \cite{Gregory2009}. Some hardware supports W-buffering but most of it doesn't and it doesn't seems to be supported in future. A workaround to this problem was presented in \cite{Dunlop2006}. It customizes projection transformation in vertex shader in way that depth buffer values have linear distribution at the end. 
\figuremacroW{edge_preservation}{Bilateral filtering results for different distances to camera}{The more distant fluid is from viewer the more smoothed its surface is. Also less edges are preserved due to z-buffer used resulting in more particles got blended into background surfaces}{0.7}

Another issue with smoothing using fixed kernel size filter is that further surfaces are more smoothed than closer surfaces (see figure~\ref{edge_preservation}). Solving that problem would require using filter with variable kernel size, depending on fragment distance from viewer. As implementing this efficiently on graphic hardware is hard and the effect is not so irritating this won't be considered any more. Section~\ref{sec:curvatureflowsmoothing} will describe another smoothing technique that doesn't suffer from this problem. 

\subsubsection{Curvature flow smoothing} \label{sec:curvatureflowsmoothing}
This technique was presented in \cite{laanSainz2009}. 

\subsection{Thickness} \label{sec:thickness}
This step is computed to determine how opaque is surface in given point. It is achieved by rendering particles as circles into depth buffer with additive blending enabled. Resulting thickness texture is then smoothed with Gaussian filter (this time regular one). In order to speed up rendering process this texture can be rendered with lower resolution and then applied with linear interpolation. 

\subsection{Rendering} \label{sec:rendering}
Last step assembles fluid surface with background image. As an input it takes fluid depth texture, thickness texture and texture with rest of the scene rendered. 
[TODO normal reconstruction, refraction, reflection, thickness]

\section{Isosurface extraction} \label{sec:isosurfaceextraction}
Classic algorithm for isosurface extraction is marching cubes (see \cite{LorensenCline1987}). It takes 3d array with scalar field values as an input and produces list of triangles. Algorithm divides space into cubes, and proceeds through scalar field taking one cube at a time (each cube consists of 8 vertices with scalar field values). Field value at each cube vertex is tested to be above or under given threshold and then appropriate triangle configuration is taken from lookup table (there are 256 triangle configurations).

This method is not suitable for extraction surface from particles. Firstly scalar field values at cube corners are not given. Secondly visiting each cube is not efficient as most of cubes doesn't contain surface. Several variations of marching cube algorithm were created to overcome those problems. The algorithm used is a multithreaded version of the one presented in \cite{RosenbergBirdwell2008}. This is a surface following version of marching cubes with a fast particle lookup cache technique. 

\subsection{Overview} \label{sec:iso_overview}
Algorithm takes as an input a list of particles positions, particle radius of influence ($R_c$), isosurface threshold and produces list of triangles. Particle radius of influence is something different than particle size (or simply particle radius - $R$), because its value is usually greater. Every particle produces a field that has non-zero values within $R_c$ from particle center. The field value in given point of space~(p) is a sum of field values generated by all particles within $R_c$. This can be represented by:
\begin{equation}
\label{eq:field_equation}
F(p) = \sum_{s~\epsilon~D_{R_c}(p)} f(dst(s, p))
\end{equation}
where $s$ is a center of particle, $D_{R_c}(p)$ is a disc with p as a center and radius $R_c$. $f(r)$ is a field function of single particle. It can be any function that is radially monotonic, continuous and has non-zero values within $R_c$ from particle center. Following function is often used as it can be computed efficiently with a few operations:
\begin{equation}
\label{eq:single_field_equation}
f(r) = \left\{ \begin{array}{ll}
{r \over \sqrt{2}R_c} )^4 - ( {r \over \sqrt{2}R_c} )^2 + 0.25 & \textrm{if $r < R_c$}\\
0 & \textrm{otherwise}\\
\end{array} \right.
\end{equation}

When $R_c > R$ neighboring particles combines into one smooth surface  (see~\ref{metaballs}). This is a common technique known as metaballs (see TODO cytat). The biggest problem here is to efficiently compute field values at corners and to traverse space in most optimal way. Those topics will be described in following subsections.
 
\figuremacroW{metaballs}{Illustration of combining particles}{SOMETHING}{1.0}

\subsection{Block subdivision} \label{sec:block_subdivision}
Space is divided into blocks containing cubes. Such a division aims for:
\begin{itemize}
\item Reduce memory consumption - smaller blocks contains less particles and cubes. 
\item Allowing for parallelize algorithm in an easy way - because processing one block is an independent of processing other blocks. 
\end{itemize}
Blocks are expanded to contain also particles that lies outside them but may influence the field values inside block. That means the block must also contain particles lying within $R_c$ distance from it - this area is called margin and have special treatment, which will be described in next section. Expanding blocks makes parallel processing very easy as there are no dependencies between them.

Block subdivision is a first step of the algorithm. For every particle containing blocks are computed. One particle can be contained in at most 8 blocks due to blocks overlapping. Each block has a list of particles lying within its boundary. 

\subsection{Block processing} \label{sec:block_processing}
Processing of block starts from building a particle lookup cache which will be described in section \ref{sec:lookupcache}. Next surface following marching cube algorithm is run inside block. The algorithm is named marching slices in \cite{RosenbergBirdwell2008}. Block containing $n x n x n$ blocks is divided on n slabs and n+1 slices (TODO rysunek). Division is made in parallel to XY plane. Each slab contains $n x n x 1$ blocks. Slices lies between slabs and serves as a cache. In addition neighboring slabs share slices that lies between them. 

\subsection{Lookup cache}\label{sec:lookupcache}
To compute field value for a given corner all particles that lies within $r_c$ from it have to be found. This task can be accomplished by building space partitioning data structure like octtree (TODO może odnośniki). However .......... . \cite{RosenbergBirdwell2008} presented another to accomplish this task. Presented method does use only integer comparison ........ .


\subsection{Multithreading}
diagram

% ---------------------------------------------------------------------------
% ----------------------- end of thesis sub-document ------------------------
% ---------------------------------------------------------------------------  % rendering techniques


% this file is called up by thesis.tex
% content in this file will be fed into the main document

%: ----------------------- introduction file header -----------------------
\chapter{Results and conclusions}

\graphicspath{{results/figures/}}

% ----------------------------------------------------------------------
%: ----------------------- introduction content ----------------------- 
% ----------------------------------------------------------------------

\section{Performance}

% table with machine configurations

Performance test were performed on two machines which configuration is presented in table~\ref{tab:test_configuration}. PhysX scene with one emitter was used. Results were rendered in two resolutions - 640~x~480 and 1280~x~960. 

%Present graphs with frame rate depending on number of particle simulated. Note that decrease in performance comes also from simulation cost, not only from rendering.

%Show second set of graphs showing time to render one frame depending on number of particles. 

%Show table showing how resolution affects performance of algorithms when number of particles is kept constant.

%Show table showing speedup of isosurface algorithm.

\begin{table}   
    \caption[Configuration of test machines]{\textbf{Configuration of test machines}}
    \centering
    \begin{tabular}{ l | p{4.5cm} | p{4.5cm} }
                                         & machine 1                                             & machine 2 \\ \hline
        CPU                           & Intel Core i7 2600k 3.4 GHz, 4 cores  & Inte Core i7 2630M 2.0 GHz, 4 cores \\ \hline
        GPU                           & Nvidia GeForce GTX 560 Titanium        &  Nvidia GeForce GT 540M \\ \hline
        Physical memory       & 8 GB                                                      & 6GB \\
    \end{tabular}
    \label{tab:test_configuration}
\end{table}

\begin{table}
   \caption[Performance comparison]{\textbf{Performance comparison} of all rendering algorithms. Screen space methods ran with 60000 particles, isosurface extraction run with 20000 particles using 3 threads.}
   \centering
   \begin{tabular} { l | r | c | r }
      Algorithm                                       &            resolution              &          machine           &          Frame (ms)      \\ \hline \hline
      Bilateral Gaussian smoothing        &           640x480                &              1                  &           10.0                    \\ \hline
      Bilateral Gaussian smoothing        &           1280x960              &              1                  &           13.8                    \\ \hline
      Bilateral Gaussian smoothing        &           640x480                &              2                  &           51.9                    \\ \hline
      Bilateral Gaussian smoothing        &           1280x960              &              2                  &           63.5                    \\ \hline \hline

      Curvature flow smoothing 100 it. &           640x480                &              1                  &           15.6                    \\ \hline
      Curvature flow smoothing 100 it. &          1280x960               &              1                  &           27.8                    \\ \hline
      Curvature flow smoothing 100 it. &           640x480                &              2                  &           73.1                    \\ \hline
      Curvature flow smoothing 100 it. &          1280x960               &              2                  &           134.1                  \\ \hline \hline

      Curvature flow smoothing 60 it.   &           640x480                &              1                  &           13.7                \\ \hline
      Curvature flow smoothing 60 it.   &          1280x960               &              1                  &           22.0                \\ \hline
      Curvature flow smoothing 60 it.   &           640x480                &              2                  &           65.7                    \\ \hline
      Curvature flow smoothing 60 it.   &          1280x960               &              2                  &           104.5               \\ \hline \hline

      Isosurface extraction                    &          640x480                &              1                  &            60.4                    \\ \hline     
      Isosurface extraction                    &          1280x960               &              1                  &           63.1                   \\ \hline      
      Isosurface extraction                    &          640x480                &              2                  &            104.8               \\ \hline
      Isosurface extraction                    &          1280x960               &              2                  &           120.5               \\ 
   \end{tabular}
\end{table}

\figuremacro{curvature_640_480}{Performance of screen space rendering with curvature flow at 640x480}{Some description}
\figuremacro{curvature_1280_960}{Performance of screen space rendering with curvature flow at 1280x960}{Some description}
\figuremacro{bilateral}{Performance of screen space rendering with bilateral Gaussian smoothing}{Some description}
\figuremacro{curvature_100it}{Performance of screen space rendering curvature flow smoothing}{100 iterations were performed for each frame.}

%CONCLUSIONS:
%- increasing resolution affects screen space rendering much more than increasing number of particles
%  - curvature flow affected much more by increasing resolution than Gaussian smoothing
%- isosurface extraction behaves in opposite way

\section{Visual appearance}

Describe what features algorithms include. Show some images. 

- isosurface extraction looks good but it's really slow, loads CPU and does not include depth.
- curvature flow produces better smoothed surface but it comes with a cost of greater complexity (computational)
- isosurface produces fluids that is more gelly like than screen space rendering


FINAL CONCLUSIONS:
- Isosurface is sufficient for small particle systems (up to ~3k - 6k particles) like blood, small fountains.
- Screen space rendering can be used to simulate bigger systems - water flooding some areas, ponds. 
   - can run bilateral gauss on lower performance hardware and curvature flow on higher.


% ----------------------------------------------------------------------





% --------------------------------------------------------------
%:                  BACK MATTER: appendices, refs,..
% --------------------------------------------------------------

% the back matter: appendix and references close the thesis


%: ----------------------- bibliography ------------------------

% The section below defines how references are listed and formatted
% The default below is 2 columns, small font, complete author names.
% Entries are also linked back to the page number in the text and to external URL if provided in the BibTex file.

% PhDbiblio-url2 = names small caps, title bold & hyperlinked, link to page 
\begin{tiny} % tiny(5) < scriptsize(7) < footnotesize(8) < small (9)

\bibliographystyle{Latex/Classes/PhDbiblio-url2} % Title is link if provided

\bibliography{9_backmatter/references} % adjust this to fit your BibTex file

\end{tiny}	

% --------------------------------------------------------------
% Various bibliography styles exit. Replace above style as desired.

% in-text refs: (1) (1; 2)
% ref list: alphabetical; author(s) in small caps; initials last name; page(s)
%\bibliographystyle{Latex/Classes/PhDbiblio-case} % title forced lower case
%\bibliographystyle{Latex/Classes/PhDbiblio-bold} % title as in bibtex but bold
%\bibliographystyle{Latex/Classes/PhDbiblio-url} % bold + www link if provided

%\bibliographystyle{Latex/Classes/jmb} % calls style file jmb.bst
% in-text refs: author (year) without brackets
% ref list: alphabetical; author(s) in normal font; last name, initials; page(s)

%\bibliographystyle{plainnat} % calls style file plainnat.bst
% in-text refs: author (year) without brackets
% (this works with package natbib)


% --------------------------------------------------------------



\end{document}
